\documentclass{book}
\usepackage{amsmath}
\usepackage{amssymb}
\usepackage{amsthm}

\theoremstyle{definition}
\newtheorem{definition}{Definition}[section]
\newtheorem{proposition}[definition]{Proposition}

\setcounter{chapter}{2}

\begin{document}

\section{Definition}

\begin{definition}
Let $V$ be a vector space over a field $F$. A \textbf{linear functional} on $V$ is a linear map $\varphi \colon V \to F$. The \textbf{dual space} of $V$, denoted $V^{*}$, is the set of all linear functionals on $V$:
\[
V^{*} = \operatorname{Hom}(V, F) = \{\, \varphi \colon V \to F \mid \varphi \text{ is linear} \,\}.
\]
\end{definition}

\begin{proposition}
$V^{*}$ is a vector space over $F$ with operations defined pointwise:
\[
(\varphi + \psi)(v) = \varphi(v) + \psi(v), \qquad
(\lambda \varphi)(v) = \lambda \cdot \varphi(v),
\]
for all $\varphi, \psi \in V^{*}$, $\lambda \in F$, and $v \in V$.
\end{proposition}

\begin{proof}
Left as an exercise.
\end{proof}

\begin{definition}
For $v \in V$ and $\varphi \in V^{*}$, we write
\[
\langle \varphi, v \rangle = \varphi(v).
\]
The map $\langle\cdot, \cdot\rangle \colon V^{*} \times V \to F$ is called the \textbf{natural pairing} between $V^{*}$ and $V$.
\end{definition}

\begin{proposition}
The natural pairing is bilinear: for all $v, w \in V$, $\varphi, \psi \in V^{*}$, and $\lambda \in F$,
\[
\begin{aligned}
\langle \varphi + \psi, v \rangle &= \langle \varphi, v \rangle + \langle \psi, v \rangle, \\
\langle \lambda \varphi, v \rangle &= \lambda \langle \varphi, v \rangle, \\
\langle \varphi, v + w \rangle &= \langle \varphi, v \rangle + \langle \varphi, w \rangle, \\
\langle \varphi, \lambda v \rangle &= \lambda \langle \varphi, v \rangle.
\end{aligned}
\]
\end{proposition}

\begin{proof}
Left as an exercise.
\end{proof}

\end{document}
